\documentclass[10pt,a4paper]{amsart}
\usepackage[margin=19mm]{geometry}
\usepackage{amsmath,amssymb,mathtools}
\usepackage[scr=rsfs,cal=euler]{mathalfa}
\usepackage{xfrac}
\usepackage[unicode,hidelinks,bookmarks=false]{hyperref}
\newcommand{\ie}{i.e.\ }
\newcommand{\cf}{cf.\ }
\newcommand{\resp}{respectively\ }
\newcommand{\Subsection}{\subsection}
\providecommand{\nobreakdash}{\nobreak\hbox{-}}
\setlength{\emergencystretch}{2em}
\multlinegap0pt
\usepackage{breqn}
\usepackage{bm}
\usepackage{amssymb}
\usepackage{xcolor}
\usepackage[shortlabels]{enumitem}
\usepackage{mathtools}
\newtheorem{lem}{Lemma}
\title{Twisted Moments of Rankin-Selberg $L$-functions in the Prime-Power Level Aspect}
\author{Fatma \c{C}i\c{c}ek, Alia Hamieh}
\date{Source: arXiv:2606.07329v1 (CC BY 4.0)}
\begin{document}
\maketitle

\noindent\emph{Source and licence.} Twisted Moments of Rankin-Selberg $L$-functions in the Prime-Power Level Aspect. Version arXiv:2606.07329v1 (CC BY 4.0), distributed by arXiv under the Creative Commons Attribution 4.0 International licence, http://creativecommons.org/licenses/by/4.0/, as recorded in the arXiv licence metadata for this exact version and shown on the abstract page https://arxiv.org/abs/2606.07329v1. Full licence notice: \texttt{licenses/arxiv-2606.07329-CC-BY-4.0.txt}.\par\smallskip

\noindent\textit{Editorial scope.} Counted unit: the article's lem lem (source0.tex lines 2287--2308). The article's complete original proof of it is delivered with it (source0.tex lines 2309--2372, one \texttt{proof} environment of 64 lines). No mathematics is written, repaired or re-derived: the statement and the proof are the article's own lines, in the article's own notation, and no retained line is substituted or re-flowed.  No further source-local context is needed: every cross-reference of every retained line resolves inside the two delivered spans.  Nothing was omitted from any retained span, so the package carries no omission record.  No retained line contains a citation, so the wrapper carries no bibliography and no bibliography entry.  No retained line contains a figure environment, an image-inclusion command or drawing code, so no image is delivered and the package declares no asset dependency.  Editorial formatting only, in addition: an AMS \texttt{amsart} preamble that restates the article's own declarations and macros used by the retained lines, namely the environments \texttt{lem} on separate counters, together with the standard packages those lines need.  The retained environments print in this document as Lemma 1 in order of appearance, whereas the article prints the same items with the numbering of its own full document.  The complete native source of the article as distributed in the arXiv e-print for this exact version is bundled unchanged as \texttt{sources/202249/source0.tex}.
\begin{lem}
For $\Re(z)>1$, we have 
 \begin{equation*}
    \sum_{n, (n, p)=1}  \frac{\lambda_g(n)  \lambda_g(aen)}{n^{z}}
    =
    \Upsilon(z, p, ae) \frac{L(z, g\otimes g)}{\zeta(2z)},
    \end{equation*}
where \[
    \Upsilon(z, p, ae)
    =  \Big( \sum_{j=0}^\infty
    \frac{\lambda_g(p^j)^2}{p^{zj}}
    \Big)^{-1}
        \prod_{q\mid ae} 
         \Big( \sum_{j=0}^\infty
    \frac{\lambda_g(q^j)^2}{q^{zj}}
    \Big)^{-1}
     \Big( \sum_{j=0}^\infty
    \frac{\lambda_g(q^j) \lambda_g(q^{\nu_{q}(ae)+j})}{q^{zj}}
    \Big)
    \]
is holomorphic and satisfies  $\left|\Upsilon(z, p, ae)\right|\ll_{\epsilon} (ae)^\epsilon$ in every fixed domain  $\Re(z)\geq\epsilon>0$. 
\end{lem}

\begin{proof}
For $\Re(z)>1$, we write   \begin{align*}
    \sum_{n, (n, p)=1}  \frac{\lambda_g(n)  \lambda_g(aen)}{n^{z}}
    &= \sum_{n}  \frac{\lambda_g(n)  \lambda_g(aen)}{n^{z}}
    \Big( \sum_{j=0}^\infty
    \frac{\lambda_g(p^j) \lambda_g(p^{\nu_p(ae)+j})}{p^{zj}}
    \Big)^{-1} \\&=  \sum_{n}  \frac{\lambda_g(n)  \lambda_g(aen)}{n^{z}}
    \Big( \sum_{j=0}^\infty
    \frac{\lambda_g(p^j)^2 }{p^{zj}}
    \Big)^{-1}   \end{align*}
since  $\nu_p(ae)=0$ provided that $(\ell, p)=1$. We further write    \[
     \sum_{n}  \frac{\lambda_g(n)  \lambda_g(aen)}{n^{z}}
     =\prod_{(q, ae)=1} 
     \Big( \sum_{j=0}^\infty
    \frac{\lambda_g(q^j)^2}{q^{zj}}
    \Big)
    \prod_{q\mid ae} 
     \Big( \sum_{j=0}^\infty
    \frac{\lambda_g(q^j) \lambda_g(q^{\nu_{q}(ae)+j})}{q^{zj}}
    \Big).
    \]
Now the definition of the $L$-function combined with this identity give
    \begin{align*}
    L(z, g\otimes g)
   & = \zeta(2z) \sum_{n=1}^\infty \frac{\lambda_g(n)^2}{n^z}
    \\
    &= \zeta(2z) 
    \prod_{q \mid ae} 
     \Big( \sum_{j=0}^\infty
    \frac{\lambda_g(q^j)^2}{q^{zj}}
    \Big)
    \prod_{(q, ae)=1} 
     \Big( \sum_{j=0}^\infty
    \frac{\lambda_g(q^j)^2}{q^{zj}}
    \Big)
    \\
    &= \zeta(2z) 
    \prod_{q \mid ae} 
     \Big( \sum_{j=0}^\infty
    \frac{\lambda_g(q^j)^2}{q^{zj}}
    \Big)
    \prod_{q\mid ae} 
     \Big( \sum_{j=0}^\infty
    \frac{\lambda_g(q^j) \lambda_g(q^{\nu_{q}(ae)+j})}{q^{zj}}
    \Big)^{-1}
     \sum_{n}  \frac{\lambda_g(n)  \lambda_g(aen)}{n^{z}} .
    \end{align*}
Thus, for $\Re(z) >1$, we have
    \begin{align*}
       & \sum_{n, (n, p)=1}  \frac{\lambda_g(n)  \lambda_g(aen)}{n^{z}}
       \\
        & = 
        \frac{L(z, g\otimes g)}{\zeta(2z) } \Big( \sum_{j=0}^\infty
    \frac{\lambda_g(p^j)^2}{p^{zj}}
    \Big)^{-1}
        \prod_{q\mid ae} 
         \Big( \sum_{j=0}^\infty
    \frac{\lambda_g(q^j)^2}{q^{zj}}
    \Big)^{-1}
     \Big( \sum_{j=0}^\infty
    \frac{\lambda_g(q^j) \lambda_g(q^{\nu_{q}(ae)+j})}{q^{zj}}
    \Big).
    \end{align*}
\end{proof}

\end{document}
